\section{Continuum theory}\label{sec:continuum}

\subsection{From Hamiltonian/Vlasov dynamics to the sheet}

Let $H:T^*\X\times[0,T]\to\R$ be a (possibly time-dependent) Hamiltonian, $X_H$ its vector field,
$\iota_{X_H}\om=\dd H$, i.e.\ $\dot x=\partial_pH$, $\dot p=-\partial_xH$, and $\Phi_{t,s}$ its flow. The Vlasov
equation for a measure $f_t$ on $T^*\X$ is, in weak form,
\begin{equation}\label{eq:vlasov}
\frac{\dd}{\dd t}\int\varphi\,\dd f_t=\int\{\varphi,H_t\}\,\dd f_t\qquad\forall\varphi\in C_c^\infty(T^*\X),\qquad
\{\varphi,H\}=\dd\varphi(X_H)=\partial_x\varphi\cdot\partial_pH-\partial_p\varphi\cdot\partial_xH .
\end{equation}

\begin{definition}[cold sheet]\label{def:sheet}
A \emph{cold sheet} is a pair $(F,\mu)$ with $F:\Q\to T^*\X$ a $C^1$ immersion and $\mu=\rho_0\,\dd q$ a positive
measure on $\Q$. It represents the measure $f=F\pf\mu$. The sheet is \emph{Lagrangian} if $F^*\om=0$ and
\emph{exact with Liouville class $c\in\R^d$} if there is a function $\Sigma$ on the universal cover $\R^d$ of $\Q$ with
\begin{equation}\label{eq:exact}
\dd\Sigma=F^*\theta,\qquad \Sigma(q+e_a)=\Sigma(q)+c_a\quad(a=1,\dots,d).
\end{equation}
We call $\Sigma$ the \emph{action} (generating function) of the sheet and $(F,\Sigma):\Q\to T^*\X\times\R$ its
\emph{Legendrian lift}; the periods $c_a=\oint_{\gamma_a}p\cdot\dd x$ along the basic cycles are its Liouville class.
\end{definition}

Since $H^1(\T^d;\R)=\R^d$, every Lagrangian sheet over $\Q=\T^d$ is exact for a unique $c$, and $\Sigma$ is unique
up to a constant.

\begin{proposition}[the sheet solves Vlasov]\label{prop:vlasov}\status{established}
Let $H_t\in C^{1,1}$ in $(x,p)$ uniformly in $t$, and $F_t=\Phi_{t,0}\circ F_0$. Then $f_t=(F_t)\pf\mu$ satisfies
\eqref{eq:vlasov}. In the self-consistent case $H_t=K(p,t)+\phi_t(x)$ with $\phi_t$ determined by the density
$\rho_t=(x_t)\pf\mu$, the same holds provided $\phi_t\in C^{1,1}(\X)$.
\end{proposition}
\begin{proof}
$\frac{\dd}{\dd t}\int\varphi\,\dd f_t=\frac{\dd}{\dd t}\int_\Q\varphi(F_t(q))\,\dd\mu(q)
=\int_\Q\dd\varphi(X_{H_t})(F_t(q))\,\dd\mu(q)=\int\{\varphi,H_t\}\,\dd f_t$. The Lipschitz hypothesis gives
existence and uniqueness of the flow. In the self-consistent case the potential is a given function of $(x,t)$
once the solution is fixed, so the same computation applies.
\end{proof}

\begin{assumption}[regular sheet solution]\label{ass:regular}
For self-consistent problems we assume that a sheet solution exists on $[0,T]$ with $\phi_t\in C^{1,1}(\X)$
uniformly in $t$. For cold data this is \emph{not} a theorem in $d\ge2$ after shell crossing (densities are
singular on caustics); in $d=1$ weak measure solutions exist \cite{ZhengMajda1994}. All structural statements below
use only that each time step applies a diffeomorphism of $T^*\X$; all error estimates are proved for externally
prescribed smooth Hamiltonians.
\end{assumption}

The configuration-space density is $\rho_t=(x_t)\pf\mu$, i.e.\ $\rho_t(x)=\sum_{q:\,x(q,t)=x}\rho_0(q)/\abs{\det D_qx(q,t)}$
away from caustics $\{\det D_qx=0\}$; the number of terms is the number of streams.

\subsection{Lagrangian structure and Poincar\'e--Cartan invariants}

\begin{proposition}[invariants of the sheet]\label{prop:invariants}\status{established}
Let $F_t=\Phi_{t,0}\circ F_0$ with $H_t$ a $C^2$ function on $T^*\X$ (single valued on the torus). Then
\begin{enumerate}[label=(\roman*),leftmargin=2.2em,itemsep=1pt]
\item $F_t^*\om=F_0^*\om$; in particular a Lagrangian sheet stays Lagrangian;
\item for every closed loop $\gamma\subset\Q$ the Poincar\'e invariant $I_t[\gamma]=\oint_{F_t\circ\gamma}\theta$ is constant in time;
\item an exact sheet stays exact with the same Liouville class $c$.
\end{enumerate}
In label coordinates $F^*\om=0$ reads
\begin{equation}\label{eq:cauchy}
\partial_ax\cdot\partial_bp-\partial_bx\cdot\partial_ap=0\qquad(1\le a<b\le d).
\end{equation}
\end{proposition}
\begin{proof}
By Cartan's formula $\Lie_{X_H}\theta=\iota_{X_H}\dd\theta+\dd\iota_{X_H}\theta=\dd(\theta(X_H)-H)$, hence
\begin{equation}\label{eq:dtheta}
\frac{\dd}{\dd t}F_t^*\theta=F_t^*\Lie_{X_{H_t}}\theta=\dd\big(\mathcal{L}_t\circ F_t\big),\qquad
\mathcal{L}_t:=p\cdot\partial_pH_t-H_t .
\end{equation}
The right-hand side is exact with a single-valued primitive because $H_t$ is a function on $T^*\T^d$. Taking
$\dd$ gives (i); integrating over a closed loop gives (ii); (iii) follows from (ii) applied to the basic cycles
together with (i).
\end{proof}

\begin{remark}[Cauchy invariants and Kelvin's theorem]\label{rem:cauchy}
With $p=a^2\dot x$, \eqref{eq:cauchy} are the Cauchy invariants of Lagrangian perturbation theory
\cite{FrischVillone2014,ZheligovskyFrisch2014,RampfVilloneFrisch2015}; in a single-stream region the sheet is the
graph of the velocity field and (ii) is Kelvin's circulation theorem. The sheet formulation extends both through
shell crossing, where the Eulerian velocity is multivalued but \eqref{eq:cauchy} persists.
\end{remark}

\begin{theorem}[action transport; Legendrian lift]\label{thm:action}\status{established mathematics; new use}
\begin{enumerate}[label=(\alph*),leftmargin=2.2em,itemsep=1pt]
\item \emph{(continuous time)} If $(F_0,\Sigma_0)$ is an exact sheet, then
$\Sigma_t(q)=\Sigma_0(q)+\int_0^t\mathcal{L}_s(F_s(q))\,\dd s$ satisfies $\dd\Sigma_t=F_t^*\theta$ for all $t$, with unchanged periods.
\item \emph{(one step of a symplectic map)} Let $\Psi$ be a diffeomorphism of $T^*\X$ with
$\Psi^*\theta-\theta=\dd A_\Psi$ for a single-valued function $A_\Psi$. If $(F,\Sigma)$ is an exact sheet then so is
$(\Psi\circ F,\ \Sigma+A_\Psi\circ F)$, with the same Liouville class.
\item \emph{(splitting integrators)} For the drift $\Psi_D(x,p)=(x+D\,\nabla K(p),\,p)$ and the kick
$\Psi_\kappa(x,p)=(x,\,p-\kappa\nabla V(x))$, with $K\in C^2(\R^d)$, $V\in C^2(\X)$ and constants $D,\kappa$,
\begin{equation}\label{eq:genfun}
A_{\Psi_D}=D\,\big(p\cdot\nabla K(p)-K(p)\big),\qquad A_{\Psi_\kappa}=-\kappa\,V(x).
\end{equation}
\end{enumerate}
Consequently a kick--drift--kick scheme with arbitrary time-dependent coefficients admits an \emph{exact}
per-vertex action update: $\Sigma\mathrel{+}=D(p\cdot\nabla K-K)$ at drifts, $\Sigma\mathrel{-}=\kappa V(x)$ at kicks.
\end{theorem}
\begin{proof}
(a) is \eqref{eq:dtheta} integrated in time. (b): $(\Psi\circ F)^*\theta=F^*\Psi^*\theta=F^*\theta+\dd(A_\Psi\circ F)$.
(c): for the drift, $p\cdot\dd(x+D\nabla K(p))-p\cdot\dd x=D\,p\cdot\dd(\nabla K(p))=D\,\dd(p\cdot\nabla K-K)$; for the
kick, $(p-\kappa\nabla V)\cdot\dd x-p\cdot\dd x=-\kappa\,\dd V$. Single-valuedness of $A_{\Psi_\kappa}$ on the torus
is the periodicity of $V$.
\end{proof}

\begin{remark}
(b) is the statement that the discrete action of a symplectic map is its generating function
\cite{MarsdenWest2001,HLW}. If $V$ were not periodic (a uniform force), $A_{\Psi_\kappa}$ would be multivalued and the
Liouville class would change by $-\kappa\oint\dd V$: the periods $c_a$ record the net impulse around the box. In
periodic cosmology the mean force vanishes and $c$ is conserved.
\end{remark}

\begin{example}[initial conditions]\label{ex:ic}
Zel'dovich data $x=q+D\nabla\chi(q)$, $p=a^2\dot D\,\nabla\chi(q)$ are exactly Lagrangian with $c=0$ and
\begin{equation}\label{eq:zeldovich}
\Sigma_0=a^2\dot D\,\big(\chi+\tfrac{D}{2}\abs{\nabla\chi}^2\big),
\end{equation}
since $p\cdot\dd x=a^2\dot D\,\nabla\chi\cdot(\dd q+D\,\dd\nabla\chi)$. Second-order data
$x=q+D_1\nabla\chi_1+D_2\nabla\chi_2$, $p=a^2(\dot D_1\nabla\chi_1+\dot D_2\nabla\chi_2)$ are \emph{not}: a direct computation gives
\begin{equation}\label{eq:2lpt}
(F^*\om)_{ab}=a^2\,(D_1\dot D_2-D_2\dot D_1)\,\big[\nabla\nabla\chi_1,\nabla\nabla\chi_2\big]_{ab},
\end{equation}
the commutator of the two Hessians, a third-order quantity in the perturbation amplitude. This is the phase-space
form of the known fact that the Cauchy invariants force a transverse displacement at third order
\cite{Buchert1994,RampfVilloneFrisch2015}. Exactly Lagrangian higher-order data with a known action can be produced
by applying exact symplectic maps (e.g.\ kick--drift steps with perturbative potentials) to Zel'dovich data and
using \cref{thm:action}(b).
\end{example}

\subsection{Reparameterisation as a gauge symmetry}

\begin{theorem}[relabelling invariance]\label{thm:gauge}\status{new derivation; elementary}
Let $(F,\Sigma,\mu)$ be an exact cold sheet and $T:\Xi\to\Q$ a $C^1$ diffeomorphism isotopic to the identity. Put
$G=F\circ T$, $\tilde\Sigma=\Sigma\circ T$, $\tilde\mu=T^*\mu:=(T^{-1})\pf\mu$, i.e.\ $\tilde\rho=(\rho_0\circ T)\det DT$. Then
\begin{enumerate}[label=(\roman*),leftmargin=2.2em,itemsep=1pt]
\item $G\pf\tilde\mu=F\pf\mu$ (the physical distribution is unchanged);
\item $G^*\om=T^*(F^*\om)=0$ and $\dd\tilde\Sigma=G^*\theta$ with the same periods $c$;
\item if $F_t$ solves Hamilton's equations and $T_t$ is any $C^1$ family of such diffeomorphisms, then
$G_t=F_t\circ T_t$, $q_t=T_t$, $\tilde\Sigma_t=\Sigma_t\circ T_t$, $\tilde\rho_t$ satisfy the arbitrary
Lagrangian--Eulerian system on $\Xi$
\begin{equation}\label{eq:ale}
\begin{aligned}
\partial_tG&=X_{H_t}(G)+(w\cdot\nabla_\xi)G,&\qquad \partial_tq&=(w\cdot\nabla_\xi)q,\\
\partial_t\tilde\Sigma&=\mathcal{L}_t(G)+(w\cdot\nabla_\xi)\tilde\Sigma,&\qquad \partial_t\tilde\rho&=\nabla_\xi\cdot(\tilde\rho\,w),
\end{aligned}
\end{equation}
with mesh velocity $w=(DT_t)^{-1}\partial_tT_t$, and conversely.
\end{enumerate}
The conclusions (i)--(ii) fail if $T$ is not injective (mass is counted with multiplicity) or does not preserve
the homology classes of the basic cycles (the periods are permuted).
\end{theorem}
\begin{proof}
(i) $G\pf T^{-1}\pf\mu=(F\circ T\circ T^{-1})\pf\mu$. (ii) Pull-back commutes with $\dd$; a map isotopic to the
identity acts trivially on $H_1$. (iii) Differentiate $G_t(\xi)=F_t(T_t(\xi))$:
$\partial_tG=X_H(G)+D_qF\,\partial_tT=X_H(G)+D_\xi G\,(DT)^{-1}\partial_tT$; the other equations are the same
computation, and the last is the continuity equation for the push-forward by $T_t^{-1}$, whose velocity is $-w$.
\end{proof}

Equation \eqref{eq:ale} shows what a remap \emph{is}: an advection in the reference coordinates of the four
fields $(G,q,\tilde\Sigma,\tilde\rho)$ by the mesh velocity. The advection is where a discretisation can damage
structure; the Hamiltonian part cannot (\cref{thm:carrier}).

\begin{theorem}[the Lagrangian condition is a momentum map]\label{thm:momentum}\status{established \cite{Donaldson1999}}
On the space $\mathcal M$ of maps $F:\Q\to T^*\X$ consider the symplectic form
$\Omega_F(\delta_1F,\delta_2F)=\int_\Q\om(\delta_1F,\delta_2F)\,\dd\mu$. The group $\Diff_\mu(\Q)$ of
$\mu$-preserving diffeomorphisms acts by $F\mapsto F\circ\eta^{-1}$, preserving $\Omega$ and every Hamiltonian of
the form $\mathcal H(F)=\int_\Q H\circ F\,\dd\mu+\tfrac12\iint\mathcal G(x(q)-x(q'))\,\dd\mu\,\dd\mu$. For a
$\mu$-divergence-free vector field $v$ on $\Q$ the function
\begin{equation}\label{eq:momentum}
J_v(F)=\int_\Q (F^*\theta)(v)\,\dd\mu
\end{equation}
generates the infinitesimal action $F\mapsto DF\cdot v$: $\dd J_v=\iota_{DF\cdot v}\Omega$. Hence $J_v$ is conserved by
the dynamics for every such $v$, and $J_v(F)=0$ for all $\mu$-divergence-free $v$ with zero mean flux if and only
if $F^*\theta$ is closed, i.e.\ $F$ is Lagrangian; the values on the harmonic (constant-flux) fields are the periods
$c_a$.
\end{theorem}
\begin{proof}
For a variation $Y$ of $F$ (a vector field along $F$), $\delta(F^*\theta)=F^*(\iota_Y\dd\theta)+\dd(\theta(Y))$, so
$\delta J_v=\int_\Q\dd\theta(Y,DF\,v)\,\dd\mu+\int_\Q\Lie_v(\theta(Y))\,\dd\mu
=\int_\Q\om(DF\,v,Y)\,\dd\mu$, the last integral vanishing because $\Lie_v\mu=0$. The characterisation of the zero
level is the Hodge--Helmholtz statement that the annihilator of the $\mu$-divergence-free fields is the space of exact one-forms.
\end{proof}

\begin{remark}[what a discretisation loses]\label{rem:noether}
\Cref{thm:momentum} is the Noether explanation of \cref{prop:invariants}: the Lagrangian property is conserved
\emph{because} the cold sheet has the relabelling symmetry of a fluid \cite{Salmon1988}. A vertex discretisation
$\mathcal M_h=(T^*\X)^N$ retains at most the permutation subgroup, so no discrete Noether theorem can deliver exact
Lagrangian structure of an interpolant. Exactness must come from a different mechanism (\cref{thm:carrier}).
The full relabelling group $\Diff(\Q)$ fibres over the space of densities with fibre $\Diff_\mu$; the
structure-relevant information of a reparameterisation is entirely in the base (the sampling density), which is
why optimal transport, the horizontal lift of this fibration \cite{Otto2001}, is the natural selector (\cref{sec:ot}).
\end{remark}

\subsection{Which perturbations of a sheet preserve the structure?}

A remap replaces the sheet by a nearby one. The following elementary statement identifies exactly which
perturbations are harmless.

\begin{proposition}[normal form of structure-preserving perturbations]\label{prop:normalform}\status{new derivation; the nonlinear version is \cite{Weinstein1971}}
Let $(F,\Sigma)$ be an exact Lagrangian immersion and $F^\varepsilon=F+\varepsilon e+O(\varepsilon^2)$ a smooth family
of immersions. Define the \emph{normal error one-form} $\beta_e\in\Omega^1(\Q)$,
\begin{equation}\label{eq:beta}
\beta_e(v):=\om(e,DF\,v)=e_x\cdot\partial_vp-e_p\cdot\partial_vx .
\end{equation}
Then
\begin{equation}\label{eq:varliouville}
\frac{\dd}{\dd\varepsilon}\Big|_{0}(F^\varepsilon)^*\theta=-\beta_e+\dd\big(\theta(e)\big),\qquad
\frac{\dd}{\dd\varepsilon}\Big|_{0}(F^\varepsilon)^*\om=\dd\beta_e .
\end{equation}
Consequently: (i) $\beta_e$ vanishes for tangential perturbations $e=DF\,v$ (relabelling) and
$e_\perp\mapsto\beta_e$ is an isomorphism from normal fields to one-forms; (ii) to first order the perturbed sheet is
Lagrangian iff $\dd\beta_e=0$, and it keeps the Liouville class iff $\beta_e$ is exact, $\beta_e=\dd\sigma$, in which case
its action is $\Sigma+\varepsilon(\theta(e)-\sigma)$; (iii) of the $d$ independent normal components of a generic
error field, structure preservation leaves exactly one scalar function $\sigma$ free. In $d=1$ every $\beta_e$ is
closed and the only obstruction is the single number $\oint_\Q\beta_e$ (the change of enclosed action).
\end{proposition}
\begin{proof}
For a family of maps with variation field $e$, $\frac{\dd}{\dd\varepsilon}F_\varepsilon^*\theta=F^*(\iota_e\dd\theta)+\dd(F^*\iota_e\theta)$
and $(\iota_e\dd\theta)(DF\,v)=-\om(e,DF\,v)$. Tangential fields give zero because $F$ is Lagrangian. Since $DF(T_q\Q)$
is a Lagrangian subspace, $\om$ identifies the quotient $T(T^*\X)/DF(T\Q)$ with $T^*\Q$, which is (i). (ii) and (iii)
are read off \eqref{eq:varliouville}.
\end{proof}

The nonlinear version is Weinstein's theorem: a neighbourhood of an embedded Lagrangian $\Lambda$ is
symplectomorphic to a neighbourhood of the zero section of $T^*\Lambda$, nearby Lagrangians are graphs of closed
one-forms, and those with the same Liouville class are graphs of $\dd\sigma$. Equivalently they are the images of
$\Lambda$ under Hamiltonian diffeomorphisms: an exact remap is indistinguishable from a small symplectic kick.
In the same spirit, every tangential (sliding) vector field on $\Lambda$ is the restriction of the Hamiltonian
vector field of a function vanishing on $\Lambda$, because $X_K$ is tangent to a Lagrangian $\Lambda$ iff
$\dd K|_{T\Lambda}=0$; reparameterisation is therefore not merely compatible with, but internal to, the symplectic category.

\begin{proposition}[cubic form]\label{prop:cubic}\status{new derivation; classical object in affine/Lagrangian geometry}
For a Lagrangian immersion $F$ define
\begin{equation}\label{eq:cubic}
C_{abc}:=\om(\partial_aF,\partial_b\partial_cF)=\partial_ax\cdot\partial_{bc}p-\partial_ap\cdot\partial_{bc}x .
\end{equation}
Then $C$ is a totally symmetric $(0,3)$ tensor field on $\Q$ (it transforms tensorially under relabelling), and it
is invariant under affine symplectic maps of phase space. For a graph $p=\nabla\phi(x)$ parameterised by $x$,
$C=D^3\phi$. In $d=1$, $C=\kappa\,\dd q^3$ with $\kappa=\om(F',F'')$, and $\abs{\kappa}^{1/3}\dd q$ is the
equi-affine arclength element of the curve.
\end{proposition}
\begin{proof}
Differentiating $\om(\partial_aF,\partial_cF)=0$ with respect to $q^b$ gives
$\om(\partial_{ab}F,\partial_cF)+\om(\partial_aF,\partial_{bc}F)=0$, i.e.\ $C_{abc}=C_{cab}$; with the obvious symmetry in
$(b,c)$ this is total symmetry. Under a change of labels the second derivatives acquire terms
$\Gamma^e_{bc}\partial_eF$, which are annihilated by $\om(\partial_aF,\cdot)$. An affine symplectic map commutes with
differentiation and preserves $\om$. The graph case is immediate.
\end{proof}

The cubic form is the lowest-order symplectic invariant of the \emph{shape} of the sheet; it will control both the
Lagrangian defect of piecewise-linear interpolation (\cref{prop:flux}) and the natural monitor function (\cref{sec:monitor}).
